\documentclass{ximera}
\title{Algebras and modules}

\begin{document}

    \begin{abstract} A reference for this material is   Jacobson, Nathan.  Basic algebra. I. Second edition. W. H. Freeman and Company, New York, 1985. xviii+499 pp. ISBN:0-7167-1480-9  \end{abstract}
    \maketitle



\begin{notation}
    From now on, all our vector spaces are over $\mathbb{R}$ unless otherwise stated.
\end{notation}

\begin{definition}[Algebras]
    Let $A$ be a real vector space. We say $A$ is an \emph{algebra} if there is an $\mathbb{R}$-bilinear associative product $A \times A \to A$.  
    
    $A$ is \emph{unital} if it has a 1.

    $A$ is commutative if $ab=ba$ for all $a,b \in A$.
\end{definition}


\begin{notation}
    From now on, all our algebras are unital unless otherwise stated.
\end{notation}



\begin{example}
    $A = \mathbb{R}$.
\end{example}

\begin{example}[Continuous functions]
    If $X$ is a topological space, then $A = C(X)$ is an algebra.
\end{example}

\begin{example}[Smooth functions]
    If $M$ is a smooth manifold, then $A = C^{\infty}(M)$ is an algebra.
\end{example}


\begin{example}[Polynomial algebra]
If $n \in \mathbb{N}$, then $A = \mathbb{R}[x_1, \ldots, x_n] $ is an algebra.
\end{example}

\begin{example}[Matrices] 
    If $n \in \mathbb{N}$, then $A = M_n (\mathbb{R})$ is an algebra. 
\end{example}

\begin{example}[Linear endomorphisms]
    If $V$ is a vector space, then
    \[  A =  End(V) : = \{ \, T : V \to V \, \vert \,  T \text{ is linear}\, \}               \]
    is an algebra. 
\end{example}


\begin{example}[Linear differential operators]
    If $n\in \mathbb{R}$, then $\mathcal{D}_n$,  the set of functions 
    \[       C^{\infty}(\mathbb{R}^n) \to C^{\infty} (\mathbb{R}^n)                  \]
    generated by the identity and  $\partial_1, \ldots , \partial_n$ under linear combinations and composition, is an algebra. 
\end{example}

\begin{example}[Tensor algebra]
If $V$ is a vector space, then 
\[  A =  \bigoplus_{j = 0 }^{\infty} V^{\otimes j}    \]
is an algebra.
\end{example}

\begin{example}[A non-unital algebra]
    Let $A = C_c (\mathbb{R})$ the set of continuous functions $f : \mathbb{R} \to \mathbb{R}$ with compact support.
\end{example}

\begin{definition}[Ideal]
    A proper subset $I \subset A$ is an \emph{ideal} if it is a sub-vector space, and 
    \[  a b , \, ba \in I  \, \text{ for all } a \in A, b \in I .     \]
    We write $I \triangleleft A$ to denote that $I$ is an ideal.

 For any subset $S \subset A$, there is a smallest ideal of $A$ containing $S$, which coinides with the intersection of all the ideals containing $S$. 
 
 \end{definition}

\begin{example}[Functions vanishing on a set] Let $A = C(X)$ and $Y \subset X$ a nonempty subset. Then
    \[     I = \{ f \in C(X) \, \vert \, f (x) = 0  \text{ for all } x \in Y  \}       \]
    is an ideal of $A$.
\end{example}




\begin{example}[Polynomials of high degree] 
 If $A = \mathbb{R} [x_1, \ldots , x_n]$ and $k \in \mathbb{N}$, then the vector space of polynomials of degree $\geq k$ is an ideal of $A$.
\end{example}

\begin{definition}[Quotient algebra]
    Let $A$ be an algebra and $I \triangleleft A$. The vector space $A/ I$ is also an algebra in a natural way.
\end{definition}

\begin{example}
    [Complex numbers] Let $A = \mathbb{R}[x]$ and $I\triangleleft A$ the ideal generated by $( x^2  + 1 ) $. Then 
    \[ A / I = \mathbb{C} . \]
\end{example}



\begin{definition}
 [Modules]   Let $A$ be an algebra. An $A$-module is an abelian group $M$ that can be multiplied by elements of $A$ in an associative way. That is, there is a function $A \times M \to M$ such that 
    \begin{itemize}
        \item $(a_1a_2 )m = a_1 (a_2m)$ for all $a_1, a_2 \in A$ and $m \in M$.
        \item $(a_1 + a_2)m = a_1 m + a_2 m$ for all $a_1, a_2 \in A$ and $m \in M$.
        \item $a(m_1 + m_2 ) = a m_1 + a m_2$ for all $a \in A$ and $m_1, m_2 \in M$.
        \item $1m = m $ for all $m \in M$.
    \end{itemize}
\end{definition}


\begin{example}[Vector spaces]
If $A = \mathbb{R}$, then the $A$-modules are exactly the vector spaces.
\end{example}

\begin{example}[The obvious one]
If $A$ is an algebra, then $A$ is an $A$-module. 
\end{example}



\begin{example}
$\mathfrak{X}(M)$ is a $C^{\infty}(M)$-module. 
\end{example}

\begin{example}
    $C^{\infty}(\mathbb{R}^n )$ is a $\mathcal{D}_n$-module. 
\end{example}


\begin{example}
    $\mathbb{R}^n$ is an $M_n(\mathbb{R})$-module. 
\end{example}


\begin{example}
   Any vector space $V$  is an $End (V) $-module. 
\end{example}


\begin{definition}[Module morphisms]
   Let $M$ and $N$ be two $A$-modules. A module morphism is a linear map $T : M \to N$ such that $aT = Ta $ for all $a \in A$. We denote by $Hom_A(M,N)$ the set of module morphisms from $M$ to $N$. If $A$ is commutative, this is again an $A$-module. 
\end{definition}

\begin{notation}
    If it is not expected to cause confusion, we write $Hom(M,N)$ instead of $Hom_A(M,N)$.
\end{notation}


\begin{definition}[Dual]
    $M^{\ast} : = Hom _A ( M , A).$
\end{definition}

\begin{definition}[Submodule]
   Let $M$ be an $A$-module and $N \subset M$. $N$ is a submodule if it is a sub-vector space and multiplying elements of $N$ by $A$ takes them within $N$. We write $N \leq _A M$ to denote the fact that $N$ is a submodule of $M$. 
\end{definition}


\begin{definition}[Span]
   Let $M$ be an $A$-module and $S \subset M$. We denote by $span_A(S) \leq _A M $ to be one (or all) of the following:
   \begin{itemize}
        \item The smallest submodule of $M$ that contains $S$.
        \item The intersections of all submodules of $M$ containing $S$.
       \item The set of linear combinations of elements of $S$ with coefficients in $A$.
   \end{itemize}
   $M$ is \emph{finitely generated} if there is a finite set $S$ with $span _A(S) = M$.
\end{definition}



\begin{proposition}[Direct sum]
If $M$ and $N$ be $A$-modules. Then the direct sum $M \oplus N$ is an $A$-module. 
\end{proposition}

\begin{proposition}[Quotient]
If $M$ is an $A$-module and $N\leq_A M$, the quotient vector space $M / N$ is also an $A$-module in a natural way. 
\end{proposition}






\end{document}